\section{Methodology}
\label{sec:method}

\subsection{Distillation Setup}
\label{sec:setup}

\paragraph{Teachers.}
To narrow the cause of differences between students down to solution style, we vary only the teacher and fix all other training conditions.
We interpolate the weights of a reasoning model ($\lambda=0$) and an instruct model ($\lambda=1$) to obtain six interpolated teachers, $\lambda \in \{0, 0.2, \dots, 1.0\}$ \citep{wortsman2022soups,ilharco2023task}.
Pair A combines R1-Distill-Qwen-1.5B and Qwen2.5-Math-1.5B-Instruct, pair B combines Qwen3-4B-Thinking-2507 and Qwen3-4B-Instruct-2507, and pair C combines R1-Distill-Llama-8B and Llama-3.1-8B-Instruct; so that element-wise mixing is meaningful, each pair consists of models branched from the same base.
Pairs B and C share their configuration with \citet{wu2026interpolation}, so the teacher-side behaviour can be checked as a replication.
Because model family and scale are confounded across pairs, we make no claim about scale effects.

\paragraph{Students.}
We use Qwen2.5-0.5B-Instruct, Llama-3.2-1B-Instruct, and OLMo-2-1B-Instruct, all in the size range where long solutions become harmful rather than helpful \citep{li2025small}.
We call each of the nine combinations of the three teacher pairs and the three students a cell.
Spreading the students over three families checks that a result is not a quirk or contamination of a single family \citep{wu2026memorization}.
OLMo-2 has a context of 4,096 tokens, so some of the longest $\lambda=0$ solutions are trained with their ends truncated.

\paragraph{Data.}
So that every $\lambda$ sees the same problems in the same number, we build the data as follows.
\begin{enumerate}
  \item We remove the problems that overlap with MATH-500 from the MATH train split and sample 4,000 problems.
  \item The teacher at each of the six $\lambda$ values solves these 4,000 problems.
  \item At each $\lambda$, we keep only solutions that are completed, answer correctly, and fit entirely within 8,192 tokens.
  \item We take only the problems that pass at all six $\lambda$ values, hold out 300 of them as an in-distribution validation set, and train on the rest (2,113, 2,261, and 1,015 problems for pairs A, B, and C).
\end{enumerate}

\paragraph{Training.}
Every student is trained with a LoRA adapter with $r=16$ and $\alpha=32$ on seven types of linear layers, with batch size 32 and learning rate $2\times10^{-4}$ for 3 epochs. The maximum length is 8,192 (4,096 for OLMo).
We use LoRA for measurement, not efficiency: the update is gathered in one pair of small matrices, so different training runs can be compared on the same footing.
For each cell and $\lambda$, we train on the full training set with three seeds. To measure the difference caused by the data alone, we also split the training set into two disjoint halves and train a separate adapter on each half, with each of the same three seeds.

\paragraph{Evaluation.}
Both teachers and students generate with $T=0.6$ and top-p 0.95; the maximum length is 8,192 tokens for teachers and 7,680 tokens for students (3,840 for OLMo).
In all experiments, the adapter of each of the three seeds solves the first 200 problems of MATH-500 twice, and the values in the tables are the mean $\pm$ standard deviation of the per-seed values.
These first 200 problems have the same difficulty and subject proportions as the full MATH-500.
Accuracy is measured as pass@1; so that truncation is not mixed in with wrong answers, we also report accuracy scored on completed responses only \citep{wang2026rollouts}.

Whether a model solves by thinking is measured by the think rate (share of responses containing \texttt{</think>}) for teachers \citep{wu2026interpolation} and by the frequency of reflective markers (\emph{wait}, \emph{recheck}, \emph{alternatively}, etc.) in the first response for students \citep{yeo2025demystifying,muennighoff2025s1}.
Because reflective markers are a proxy, we also report token counts and the rate of hitting the generation limit \citep{chen2026molecular,goyal2026cap}, and for conditions that repeat the same sentence we count markers in the first 300 words.
The teacher's think rate and reflective markers are counted, in the same way as for students, over all of the teacher's solutions to the 4,000 problems.

To check whether a result arises only from the frequency of reflective words, we look again through internal features of the Llama student, for which a public SAE (sparse autoencoder) exists,\footnote{\url{https://huggingface.co/EleutherAI/sae-Llama-3.2-1B-131k}} and discard features that turn on or off merely when reflective words are inserted or deleted \citep{ma2026sae}.
We select features that, when the untrained Llama student reads the solutions of the $\lambda=0$ and $\lambda=1$ teachers, fire much more strongly on the $\lambda=0$ side, and keep only those that pass two falsification checks, inserting and deleting reflective words; we call these thinking features below.
Because we apply an SAE trained on the base model to the untrained student and to the students with adapters, we also check its reconstruction error.
We measure the mean activation of the selected features as each adapter reads its own first responses, normalized so that the activation on the $\lambda=0$ teacher solutions is 1.

To check whether the results hold outside MATH, we evaluate the adapters trained on MATH, unchanged and with the same settings, on the first 200 problems of GSM8K \citep{cobbe2021training} and GPQA-Diamond \citep{rein2023gpqa} (for GPQA-Diamond, all 198 problems).
To gauge how much distillation changed the students, we also evaluate the untrained students with the same settings; since they have no adapter, we vary three generation seeds instead of adapter seeds and report the mean $\pm$ standard deviation.

\subsection{Subspace Analysis}
\label{sec:subspace}

\paragraph{Decomposition.}
\begin{equation}
\Delta W = \tfrac{\alpha}{r} BA = U\Sigma Q^\top, \qquad U = [\,u_1, \dots, u_r\,]
\label{eq:1}
\end{equation}
\begin{equation}
\phi(\Delta W_i, \Delta W_j) = \lVert U_i^\top U_j \rVert_F^2 / r
\label{eq:2}
\end{equation}
We read the directions along which the adapter at each $\lambda$ updates the weights from the top $r$ left singular vectors $U$ of each layer's update (Eq.~\eqref{eq:1}).
Whether two adapters use the same directions is read from the overlap $\phi$ of their $U$ (Eq.~\eqref{eq:2}); it is 1 for identical directions and $r/d$ between random subspaces, where $d$ is the output dimension of the layer \citep{hu2021lora}.
Values for individual layers and modules are averaged with weights $\lVert \Delta W \rVert_F$.

\paragraph{Comparison.}
\begin{equation}
R(\lambda) = \dfrac{\tfrac{1}{2}\left[\,\phi\!\left(\Delta W_a^{0},\, \Delta W_b^{\lambda}\right) + \phi\!\left(\Delta W_b^{0},\, \Delta W_a^{\lambda}\right)\right]}{\phi\!\left(\Delta W_a^{\lambda},\, \Delta W_b^{\lambda}\right)}
\label{eq:3}
\end{equation}
Whether a change of direction is due to the teacher is read from $R$ (Eq.~\eqref{eq:3}), the overlap between adapters with different teachers divided by the overlap between adapters that differ only in data.
Unlike prior work, which did not compare overlaps with chance level or with data differences, we decide only by this ratio \citep{dieudonne2026adapter}.
Because the data halves $a$, $b$ and the seed enter the numerator and the denominator equally, only the teacher effect remains; $R=1$ means the teacher effect is at the level of the data difference, and $R$ near 0 means the adapter sharply changes the directions along which it updates the weights.
We do not use the overlap between adapters that differ only in seed as the baseline, because adapters trained with the same seed share their LoRA initialization, which inflates their overlap; we also inspect the numerator and the denominator separately.

Whether the adapter sharply changes its directions at the $\lambda$ where the teacher stops thinking is tested by a permutation test on the correlation, over $\lambda$ intervals, between the drop in $R$ and the drop in the teacher's think rate.
We do not choose a single switch $\lambda$ in advance, because the answer would depend on the criterion and the choice would be post hoc.
Since the three students of a pair share a teacher, the statistic is the mean correlation of the three students in each pair, and an exact $p$ is obtained from all ways of permuting the teacher order of each of the three pairs ($120^3$).
On the teacher side, we repeat the same test with the drop in reflective markers in the solutions instead of the think rate.

\paragraph{Extraction.}
\begin{equation}
P = \frac{1}{n}\sum_{t=1}^{n} U_t U_t^\top,\qquad V \leftarrow \{\,v : Pv = \mu v,\ \mu > 0.5\,\}
\label{eq:4}
\end{equation}
To find the directions that the pre-switch adapters all use, we take as the switch interval the first interval in which the teacher's think rate drops by more than half, overlap the $U$ of the adapters before it, and keep only the directions $V$ that more than half of them share (Eq.~\eqref{eq:4}) \citep{gerasimov2026unstable}. We call $V$ the shared region below.
The degree to which each adapter uses $V$ is measured as $\lVert U_t^\top V \rVert_F^2 / k$, where $k$ is the dimension of $V$, and compared with the random level $r/d$; both scoring and editing use a $V$ built without that adapter.

To see whether the shared region is an artifact of the design that interpolates teachers, we also build a region from only two $\lambda=0$ adapters trained on the same data with different seeds (the two-seed region below).
With two adapters, the eigenvalue of a paired direction is $(1+\cos\theta)/2$, so the 0.5 threshold would keep almost every direction; we therefore keep only directions that overlap within $60^\circ$ ($\mu > 0.75$).
The overlap between the two regions is measured as the fraction of the shared region's dimensions covered by the two-seed region $V_{\text{seed}}$, $\lVert V_{\text{seed}}^\top V \rVert_F^2 / k$.

\subsection{Subspace Interventions}
\label{sec:interventions}

\paragraph{Removal.}
\begin{equation}
\Delta W' = (I - VV^\top)\,\Delta W
\label{eq:5}
\end{equation}
To test whether the shared region is needed for solving by thinking, we remove only the $V$ component from the $\lambda=0$ adapter (Eq.~\eqref{eq:5}) \citep{arditi2024refusal}.
We implement this as $B \leftarrow B - V(V^\top B)$, so $A$ and the rank are unchanged.

\paragraph{Grafting.}
\begin{equation}
\Delta W'(\lambda{=}1) = \Delta W(\lambda{=}1) + \beta\, VV^\top \Delta W(\lambda{=}0)
\label{eq:6}
\end{equation}
To test whether adding only the shared region is enough to make the student solve by thinking, we add only the $V$ component of the $\lambda=0$ adapter to the $\lambda=1$ adapter, which gives only short answers (Eq.~\eqref{eq:6}) \citep{arditi2024refusal}.
We use $\beta=1$ and $2$; at 4 and 8 every condition repeats the same sentence, so we exclude them from the comparison. As controls, we add the component outside $V$ and the whole update with the same energy.

\paragraph{Amplification.}
\begin{equation}
\Delta W' = \lVert \Delta W \rVert_F \cdot \dfrac{\Delta W + \gamma\, VV^\top \Delta W}{\lVert \Delta W + \gamma\, VV^\top \Delta W \rVert_F}
\label{eq:7}
\end{equation}
To test whether the shared region is also a dial that sets the amount of thinking, we increase only the weight on $V$ while keeping the norm fixed (Eq.~\eqref{eq:7}) \citep{ma2025cotvalve,galichin2026covered,sheng2025planning}.
We use $\gamma=1$ and $4$, and also examine the case in which the whole update is moved into $V$.

\paragraph{Controls.}
To test whether an effect is due to the region $V$ itself, we compare with conditions that remove the same energy from the highest-energy directions orthogonal to $V$ \citep{lion2025polar,fu2026scalpel}, that shrink the whole update by the same amount, or that remove the same energy starting from the largest singular components \citep{lee2025goldswitch}.
Whether a gain in accuracy is merely due to shorter responses is checked with a condition that cuts the original adapter at the same length \citep{muennighoff2025s1}, and whether it is a side effect of the edit itself is checked with a condition that applies the same edit to the $\lambda=1$ adapter, which barely uses the shared region.

To see how cleanly the two-seed region, found without interpolating teachers, turns thinking off, we compare it with removing the largest components by the Gold-Switch criterion \citep{lee2025goldswitch}, inserting a sentence saying that thinking is finished before the answer \citep{ma2025nothinking}, cutting at the same length and having the student write the answer \citep{muennighoff2025s1}, merging half-and-half with the short-solution adapter \citep{wu2025merging}, and the untrained student.
Since the students have not learned the \texttt{</think>} tag, we insert only the sentence, without the tag.
